% TOP_PROBLEM: {"id": 2792, "title": "Kirby Problem 2.44", "category_id": 7, "category": {"id": 7, "name": "topology", "display_name": "Topology", "description": "Properties preserved under continuous deformations.", "slug": "topology", "order_index": 7, "created_at": "2026-07-31T15:26:25.671Z"}, "status": "open", "problem_number": "KP-2.44", "set_id": 11, "statusQualification": "Makes completeness and the distinction between a complete interior metric and a compact geodesic-boundary model explicit."}
% FIRST_POSTED: 2026-10-01
% ENTRY_KIND: statement_only
% UnsolvedMath ID 2792; source catalogue code KP-2.44
% !TeX program = lualatex
% Definition and sources only; this file is not an LLM solution attempt.
\documentclass[11pt,a4paper]{article}
\usepackage[margin=25mm]{geometry}
\usepackage{fontspec}
\setmainfont{lmroman10-regular.otf}[BoldFont=lmroman10-bold.otf,
 ItalicFont=lmroman10-italic.otf,BoldItalicFont=lmroman10-bolditalic.otf]
\usepackage{amsmath,amssymb,mathtools}
\usepackage{unicode-math}
\setmathfont{latinmodern-math.otf}
\AtBeginDocument{\let\setminus\smallsetminus}
\usepackage{xurl,hyperref}
\hypersetup{colorlinks=true,urlcolor=blue}
\setcounter{secnumdepth}{0}
\setlength{\parindent}{0pt}
\setlength{\parskip}{5pt}
\setlength{\emergencystretch}{3em}
\begin{document}
\section*{Kirby Problem 2.44}\label{2792}
{\small UnsolvedMath ID 2792 · KP-2.44 · Topology · Kirby's Problems in Low-Dimensional Topology}
\subsection{Definitions and mathematical statement}
Let $S$ be a connected orientable surface of infinite type: it is not homeomorphic to a compact surface with finitely many points removed. Work without boundary, or use interiors when a boundary convention is needed. Given a homeomorphism $f:S\to S$, form the mapping torus
\[
M_f=(S\times[0,1])/\bigl((x,1)\sim(f(x),0)\bigr).
\]
It is a three-manifold fibering over the circle, with fiber $S$. No compactness or finite-volume assumption is imposed on $M_f$.

Characterize the homeomorphisms $f$ for which $M_f$ admits a complete Riemannian metric of constant sectional curvature $-1$. The metric need not be unique, and completeness is part of the hyperbolic-structure requirement. Criteria may depend on the action of $f$ on curves, subsurfaces, and the ends of $S$; finite-type pseudo-Anosov terminology alone is not a definition of the desired infinite-type class.

For mapping tori that are hyperbolic, ask also whether $M_f$ is homeomorphic to the interior of a compact hyperbolic three-manifold with totally geodesic boundary. The latter means a compact manifold carrying a curvature $-1$ metric whose boundary components have zero second fundamental form. This is a topological identification with that interior, not a requirement that its incomplete restricted metric equal the chosen complete metric on $M_f$.

The second assertion is a further finiteness and geometric condition, and must be distinguished from the first existence question.
\subsection{Short English statement}
Which infinite-type surface maps have hyperbolic mapping tori, and when do those tori admit compact cores with totally geodesic boundary?
\subsection{Sources}
\begin{itemize}
\item[S1] ULAM.AI and the UnsolvedMath Contributors, supplied problems.json, record 2792 (KP-2.44), Kirby's Problems in Low-Dimensional Topology. Catalogue identity and source transcription; CC BY 4.0.
\url{https://huggingface.co/datasets/ulamai/UnsolvedMath}.
\item[S2] K3: A New Problem List in Low-Dimensional Topology, Problem 2.44. Expanded formulation of the supplied catalogue statement.
\url{https://math.berkeley.edu/sites/default/files/surv-295-ruberman-watermarked-author-pdf.pdf}.
\end{itemize}
\end{document}
